\documentclass[11pt,a4paper]{article}
\usepackage[T1]{fontenc}
\usepackage[utf8]{inputenc}
\usepackage[english]{babel}
\usepackage{amsmath,amssymb}
\usepackage[margin=2cm]{geometry}
\setlength{\parindent}{0pt}
\setlength{\emergencystretch}{3em}
\allowdisplaybreaks

\title{About \texttt{metric\_identities.tex}}
\author{}
\date{}
\begin{document}
\maketitle
This note explains what the test data in \texttt{metric\_identities.tex} is used for. That file is input for the engine in \texttt{main.c}; it is read by the test \texttt{test\_index\_identities.c} in this folder.

\begin{enumerate}
\item \textbf{Purpose.} The file holds 35 true identities of tensor calculus in which upper and lower indices are distinguished. No metric is ever assumed: each statement declares a diagonal metric with rational entries among its hypotheses, and the value of the Levi-Civita symbol when it appears. The test asks the engine to prove each one and to produce a certificate that the independent checker accepts.

\item \textbf{Where the statements come from.} The identities that the sources state in space-time, with the metric they use written out: $\varepsilon^{\mu\nu\rho\sigma}\varepsilon_{\mu\nu\rho\sigma} = -24$ and $\varepsilon^{\mu\nu\rho\sigma}\varepsilon_{\mu\nu\rho\tau} = -6\delta^{\sigma}_{\tau}$ (signature and sign of $\varepsilon$ as in Jackson, used by Baird and by the Maryland solution), $g^{\mu\nu} g_{\mu\nu} = 4$ (Flogen), $g^{ik} g_{km} = \delta^{i}_{m}$ (Carati, Barbieri), the four ways of writing a trace (Roma, exercise 3), and equation (9) of L.~Stein's problem set with the Euclidean metric declared. They are written with Latin indices running from 1 to $d$, because Greek indices and indices from 0 to 3 are not read yet. Every statement was confirmed by the oracle of \texttt{tools/fuzz/fuzz\_tensors.py} in its \texttt{data} mode, which raises and lowers indices in its own way.

\item \textbf{What the engine needs in order to pass.} To read the hypotheses, to check that each summed index appears once above and once below and each free index in the same position everywhere, to unfold each raised or lowered index into a contraction with the metric or its inverse, and to check that the unfolded identity has, at every point, the values of the written one.

\item \textbf{Format expected by the test.} One \texttt{\textbackslash item} per statement, written as \texttt{\$\textbackslash displaystyle ...\$}. The hypotheses come first, joined by \texttt{\textbackslash land}, then \texttt{\textbackslash implies} and the identity: the metric as $g_{ij} = \operatorname{diag}(c_1, \dots, c_d)$, and the value of $\varepsilon$ as $\varepsilon_{12 \dots d} = 1$ or $\varepsilon^{12 \dots d} = 1$.

\item \textbf{How to run it.} From the top folder of the repository:
\begin{verbatim}
cc -std=c11 -O2 -Wall -Wextra -Werror -o /tmp/test_index_identities \
    tests/tensor_calculus/test_index_identities.c && /tmp/test_index_identities
\end{verbatim}

\item \textbf{Expected result.} The lines \texttt{PASS: 35 of 35 statements proved, each with a written proof}, \texttt{PASS: 35 of 35 statements proved, 35 with a certificate the checker accepts} and \texttt{PASS: 35 of 35 altered certificates rejected for the right reason}.
\end{enumerate}
\end{document}
